\section{Related Work}
\label{sec:related}

\subsection{TextVQA and text-grounded visual question answering}

TextVQA was introduced to measure a model's ability to answer questions whose solution depends on reading text embedded in real-world images \citep{singh2019textvqa}. Unlike generic VQA, the benchmark frequently requires exact recognition of sparse text, reasoning over OCR tokens, and answer generation under strict string-matching evaluation. This makes the problem particularly sensitive to small prompt decisions, OCR formatting, and answer normalization. For an empirical paper, the implication is straightforward: a convincing TextVQA study must treat text handling as part of the research question rather than as an implementation footnote.

That requirement shapes the present study in two ways. First, OCR-aware prompting is treated as a structured experimental variable throughout zero-shot evaluation rather than as a single fixed preprocessing choice. Second, OCR remains part of the fine-tuning analysis through explicit on/off follow-up runs. This design choice is motivated by the benchmark itself: if scene text is central to the task, then conclusions about backbone quality or adaptation quality are incomplete unless OCR usage is interrogated directly.

\subsection{Open-source vision-language models}

Open-source VLMs have advanced rapidly through increasingly effective combinations of pretrained vision encoders, pretrained language models, and instruction-tuning pipelines. BLIP-2 demonstrated that strong multimodal transfer can be achieved by bootstrapping from frozen unimodal components rather than retraining the full stack from scratch \citep{li2023blip2}. LLaVA showed that visual instruction tuning can substantially improve general-purpose multimodal behavior with a comparatively lightweight alignment stage \citep{liu2023visual}. More recent model families, including Qwen2.5-VL \citep{qwen2025qwen25vl} and InternVL-style systems, have pushed open-source multimodal performance far beyond early proof-of-concept systems, making it realistic to benchmark several families within one study.

However, strong checkpoints alone do not guarantee a rigorous comparison. Public model reports typically emphasize aggregate benchmark leaderboards, while a project like ours must instead decide how to compare models under a fixed, transparent protocol. In particular, an academic benchmark on TextVQA must choose how much prompt breadth to allow, which split should be used for screening, and when it is appropriate to spend the official validation budget. Our funnel design can be understood as a methodological answer to that problem: broad internal-development screening is used to expose prompt and backbone sensitivity, while the official validation split is deliberately held back for a much smaller set of finalists and final tuned-model checks.

\subsection{Parameter-efficient adaptation}

The fine-tuning stage in this project is built around Low-Rank Adaptation (LoRA) \citep{hu2022lora}, which remains the most practical parameter-efficient adaptation method for a study that aims to compare several adaptation settings rather than optimize a single maximal model. LoRA is especially appropriate here for three reasons. First, it makes a multi-run ablation feasible within a limited compute budget. Second, it allows adaptation choices such as target-module strategy and rank to be varied systematically. Third, it keeps the training question close to the paper's actual scope: we are not trying to establish a new state of the art on TextVQA, but rather to understand how far a strong open-source backbone can be pushed under a controlled and reproducible adaptation budget.

Most importantly, we treat LoRA not as a one-off fine-tuning step but as a structured study. The core matrix compares attention-only and all-linear target strategies, rank $16$ versus $32$, and seed stability. The follow-up runs then ask higher-level questions that the core matrix alone cannot answer: whether OCR conditioning materially changes adapted performance, and whether additional data continues to improve the selected family. This winner-conditioned design is more informative than a single tuned checkpoint and more defensible than ablations attached to a family that never actually won the initial training comparison.
